\section{Introduction}
\label{sec:intro}

The convergence of Open RAN (O-RAN) principles with emerging 6G
technologies is reshaping how wireless networks are designed, deployed,
and optimised~\cite{oran_alliance2022,6g_survey2023}.
By disaggregating the radio access network into open, interoperable
components and exposing standardised interfaces (E2, A1, O1), O-RAN
creates a programmable substrate on which data-driven control loops can
be built and evaluated.

NS-3 has long served as the de-facto open-source network simulator for
protocol-level research~\cite{ns3_main}; however, an integrated
framework that combines NS-3 with an O-RAN stack and 6G features---THz
channels, ultra-massive MIMO, and MEC---has been lacking.
This paper fills that gap.

Our contributions are:
\begin{enumerate}
  \item A modular NS-3 O-RAN 6G integration framework that implements the
        near-RT RIC, E2 interface stubs, and xApp abstractions.
  \item An RL-based handover optimisation xApp and a federated learning
        (FL) sub-framework for distributed KPI prediction.
  \item A digital-twin module for proactive KPI forecasting and
        control-loop validation.
  \item Reproducible simulation campaigns with automated result export
        (CSV/JSON) and LaTeX paper asset generation.
\end{enumerate}

The remainder of the paper is organised as follows.
Section~\ref{sec:related} surveys related work.
Section~\ref{sec:architecture} describes the system architecture.
Section~\ref{sec:algorithms} details the algorithms.
Section~\ref{sec:setup} presents the experimental setup.
Section~\ref{sec:results} reports and analyses results.
Section~\ref{sec:discussion} discusses findings and limitations.
Section~\ref{sec:conclusion} concludes.
